\chapter{Conclusioni}\label{chapter:conclusione}

\section{Le risposte alla domanda}\label{sec:risposte}

La domanda (sezione~\ref{sec:domanda}) si articolava in quattro parti. Le prime
tre hanno una risposta misurata; la quarta aspetta i giudizi del
capitolo~\ref{chapter:confronto}.

\paragraph{I difetti di Koskidex.} Dei quattro difetti sospettati leggendo il
codice, il più grande non era fra quelli sospettati. Il recupero congiuntivo
lasciava a vuoto 290 query di SciFact su 300: correggerlo porta il nDCG@10 da
0,0246 a 0,4936, e solo dopo ha senso parlare di punteggio. BM25 lo porta a
0,6197, e la correzione di un difetto dentro BM25, le espansioni pesate con la
frequenza del termine trovato, a 0,6694; con le stopword Koskidex arriva a
0,6757, il 99,5\% del riferimento di Lucene, e su NFCorpus al 95\%. Sui vettori, il
re-ranking vale due punti con il recupero disgiuntivo; con quello congiuntivo di
Documentale non vale niente, e l'unione dei candidati vale +0,62 su SciFact. La
costante della fusione era otto volte troppo bassa per le query in lingua naturale
e giusta, per caso, per quelle identificative.

\paragraph{I difetti di Documentale.} La ricerca di produzione ha tre difetti che
colpiscono le ricerche più ordinarie: un numero d'atto trova gli atti con un numero
a un refuso, un articolo eliso rende introvabile una parola in quattro atti su
dieci, e numero e comune insieme non trovano quasi mai l'atto, perché devono stare
nello stesso campo. Con Koskidex innestato e tre impostazioni, raccolte in un
profilo, l'atto giusto passa dall'1,3\% al 91,7\% di prime posizioni sulle
ricerche per numero e comune, e le parole introvabili dal 40,3\% allo 0,4\% degli
atti. Nessuno dei tre è un errore di programmazione: sono impostazioni
ragionevoli di un motore generico, che si rompono su proprietà precise dei
documenti amministrativi italiani.

\paragraph{Una configurazione per tutto.} Non c'è. Il recupero congiuntivo è giusto
per le ricerche identificative e disastroso per le frasi; le vie di mezzo di
Lucene ed Elasticsearch non reggono tutte e due; il peso giusto del vettore
dipende dal tipo di query. Sceglierlo con una regola di una riga, \emph{recupero
congiuntivo e vettore leggero se la query ha una cifra e il congiuntivo trova
qualcosa}, vale +0,023 di nDCG@10 medio sulla migliore configurazione fissa. Se la
regola riconosca il tipo di query o la collezione, su questi dati non si può dire.

\paragraph{Il confronto con Elasticsearch.} Sugli stessi insiemi di risultati,
Koskidex risponde in una mediana di 2,7~ms contro 22, su due motori che non girano
alla pari. Quanto valgano in pertinenza, Koskidex corretto ed Elasticsearch di
produzione, sulle ricerche in lingua naturale, lo dirà il
capitolo~\ref{chapter:confronto}.

\section{Limiti}\label{sec:limiti}

\begin{itemize}
  \item \textbf{Query fabbricate.} Le known-item automatiche hanno una sola forma,
        \texttt{<numero> <comune>}, e le 24 query del confronto le ha scritte
        l'autore. Misurano bene quello che misurano, e niente di più. Le query
        scritte da persone servono proprio a questo, e al momento di scrivere non
        ci sono ancora.
  \item \textbf{Un solo corpus di dominio.} Gli albi pretori hanno la forma di un
        documentale, ma sono gli atti di comuni: altri archivi, con contratti o
        corrispondenza, hanno altri vocabolari e altre ricerche. Il testo intero
        c'è solo per 563 atti su 10.018.
  \item \textbf{Tipi di query separati per collezione.} Le query identificative
        stanno tutte sugli albi, quelle in lingua naturale sulle collezioni
        pubbliche: una regola che riconosce la collezione e una che riconosce il
        tipo di query danno gli stessi numeri. È il limite più serio della scelta
        per tipo di query, e la ragione per cui resta fuori dal motore.
  \item \textbf{Pochi annotatori.} I giudizi del pool vengono da un annotatore, con
        un secondo su una parte; il kappa dirà quanto dipendono dalla persona.
  \item \textbf{Il richiamo del pool è un limite superiore.} I documenti pertinenti
        che nessuna delle dodici configurazioni ha messo fra i primi dieci non
        sono giudicati, e contano come non pertinenti.
  \item \textbf{I tempi.} Koskidex ed Elasticsearch non girano alla pari: uno
        nativo, l'altro in una macchina virtuale e misurato attraverso una
        chiamata HTTP. I tempi dicono l'ordine di grandezza, non un rapporto.
  \item \textbf{Un solo modello di embedding.} \texttt{bge-m3} in locale è una
        scelta, non un confronto: con un altro modello i pesi calibrati del
        capitolo~\ref{chapter:ibrido} potrebbero cambiare.
\end{itemize}

\section{Sviluppi futuri}\label{sec:sviluppi}

\begin{itemize}
  \item \textbf{Le known-item umane}, e con loro tre misure che le aspettano: la
        scelta per tipo di query su query dei due tipi nella stessa collezione,
        stopword e stemmer italiani, e la scheda contro il testo intero.
  \item \textbf{Un secondo stadio di riordino.} Un modello neurale che riordina i
        primi cento candidati (\emph{cross-encoder}) direbbe quanto vale
        riordinare, dato il richiamo del primo stadio: il re-ranking funziona se è
        progettato come stadio, con pochi candidati e un modello preciso. Non
        dentro Koskidex, che resterebbe senza dipendenze, ma come servizio a parte.
  \item \textbf{La scansione dei vettori.} A 10.000 documenti una scansione esatta
        costa 12-15~ms per passata; più accumulatori nel prodotto scalare la
        accorcerebbero, a prezzo di punteggi diversi nelle ultime cifre, e oltre
        qualche centinaio di migliaia di documenti servirebbe un indice
        approssimato.
  \item \textbf{Koskidex in produzione.} La tesi misura il recupero; un motore in
        produzione chiede anche un'istantanea che non riscriva l'indice intero,
        l'aggiornamento parziale di un documento, e un indice che non stia tutto
        in memoria.
\end{itemize}
